\documentclass[aspectratio=169,11pt]{beamer}

\usetheme{Madrid}
\usecolortheme{default}
\usefonttheme{professionalfonts}
\setbeamertemplate{navigation symbols}{}
\setbeamertemplate{footline}[frame number]

\usepackage{amsmath, amssymb, booktabs, array}

\title{LGBTQ+ Identities, Work, Discrimination, And Labor-Market Institutions}
\subtitle{Gender Study --- Optional Long Lecture After Week 9}
\author{Teaching deck draft}
\date{}

\begin{document}

\begin{frame}
  \titlepage
\end{frame}

\begin{frame}{Central question}
\begin{itemize}
    \item How do sexual orientation, gender identity, disclosure, workplace climate, legal recognition, and benefits shape labor-market outcomes?
    \item Why does LGBTQ+ labor research belong inside labor economics rather than outside the field?
    \item Which margins do current designs identify, and which margins remain hidden?
\end{itemize}
\vspace{0.25cm}
\[
\text{identity object} \rightarrow \text{observed labor margin} \rightarrow
\text{identification} \rightarrow \text{welfare interpretation}.
\]
\end{frame}

\begin{frame}{Why this is labor economics}
\begin{itemize}
    \item The outcomes are labor outcomes: employment, hours, wages, earnings, occupations, firms, benefits, retention, and job quality.
    \item The mechanisms are familiar: screening, sorting, bargaining, firm treatment, policy incidence, household specialization, and amenities.
    \item The measurement problem is sharper: orientation, identity, legal category, perceived identity, and disclosure can differ.
    \item The welfare object is broader than pay: climate, concealment, harassment, safety, dignity, and benefit access matter.
\end{itemize}
\end{frame}

\begin{frame}{Bridge from Weeks 1--9}
\scriptsize
\begin{tabular}{p{0.18\linewidth} p{0.33\linewidth} p{0.38\linewidth}}
\toprule
Course block & Earlier object & Week 10 LGBTQ+ translation \\
\midrule
Week 1 & gender measurement & orientation, gender identity, disclosure, perception \\
Weeks 2--3 & household allocation, human capital, sorting & family recognition, expectations, occupation choice \\
Week 4 & hiring, promotion, firms & audits, callbacks, firm-specific barriers \\
Week 5 & norms and institutions & stigma, disclosure, coworker and manager climate \\
Week 6 & policy incidence & protection, marriage, benefits, insurance \\
Week 7 & safety and hidden harms & harassment, concealment, dignity, retention \\
Week 8 & comparative mechanisms & law, enforcement, benefits, public/private sectors \\
Week 9 & data and methods & small cells, identity linkage, climate measures \\
\bottomrule
\end{tabular}
\end{frame}

\begin{frame}{Measurement is the estimand}
\scriptsize
\begin{tabular}{p{0.27\linewidth} p{0.31\linewidth} p{0.30\linewidth}}
\toprule
Measured object & Good for & Main limit \\
\midrule
Same-sex couple & household labor supply, benefits, specialization & misses unpartnered people and direct identity \\
Direct orientation question & self-identified orientation and outcomes & disclosure bias, sample size \\
Gender identity module & trans and nonbinary outcomes & small cells, privacy, wording changes \\
Administrative marker change & linked earnings dynamics & selected high-fidelity subgroup \\
Resume identity signal & employer response at screening & visible signal, not full population \\
Climate survey & hidden harms and support & hard to link to careers \\
\bottomrule
\end{tabular}
\end{frame}

\begin{frame}{Key empirical distinctions}
\begin{columns}[T,onlytextwidth]
\begin{column}{0.48\linewidth}
\textbf{Identity object}
\begin{itemize}
    \item sexual orientation
    \item gender identity
    \item legal category
    \item perceived identity
    \item disclosure or visibility
\end{itemize}
\end{column}
\begin{column}{0.48\linewidth}
\textbf{Labor margin}
\begin{itemize}
    \item applications and callbacks
    \item employment and earnings
    \item benefits and insurance
    \item promotion and retention
    \item climate, harassment, welfare
\end{itemize}
\end{column}
\end{columns}
\end{frame}

\begin{frame}{Employment, earnings, occupations, benefits}
\begin{itemize}
    \item Reviews show heterogeneous earnings and employment patterns across gay men, lesbians, bisexual workers, trans workers, nonbinary workers, couples, and settings.
    \item Administrative identity-linkage designs can observe earnings dynamics, but usually for selected measured groups.
    \item Occupation and firm gaps may reflect search, anticipated treatment, benefits, safety, climate, or prior exclusion.
    \item Employer-sponsored benefits are compensation: partner benefits, spousal coverage, dependent coverage, and leave can change welfare even without a wage change.
\end{itemize}
\end{frame}

\begin{frame}{Hiring discrimination evidence}
\scriptsize
\begin{tabular}{p{0.25\linewidth} p{0.30\linewidth} p{0.32\linewidth}}
\toprule
Design family & Identifying variation & Observed margin \\
\midrule
Tilcsik-style audit & randomized gay identity signal on resumes & callback response \\
Sexual orientation and presentation audits & randomized application content & invitation or callback \\
Transgender hiring audit & randomized trans versus cis signal & employer response \\
Large-employer audit platform & randomized signals across many firms & callback plus firm heterogeneity \\
\bottomrule
\end{tabular}
\vspace{0.25cm}
\begin{itemize}
    \item These designs are clean for screening-stage discrimination.
    \item They do not identify post-hire wages, promotion, retention, disclosure costs, or workplace climate.
\end{itemize}
\end{frame}

\begin{frame}{Workplace climate, disclosure, harassment}
\begin{itemize}
    \item Disclosure is not the same object as identity. It can be a choice under constraint or a response to expected treatment.
    \item Climate affects retention, search, promotion, voice, safety, wellbeing, and willingness to use benefits.
    \item Survey and list-experiment designs can measure stigma or support when direct reporting is sensitive.
    \item A zero wage gap does not imply zero harm if workers avoid firms, conceal identity, or bear hidden costs.
\end{itemize}
\vspace{0.15cm}
\[
V = w - \phi h - \rho r - \psi H + B + A
\]
\end{frame}

\begin{frame}{Law, family recognition, benefit design}
\begin{columns}[T,onlytextwidth]
\begin{column}{0.32\linewidth}
\textbf{Protection}
\begin{itemize}
    \item employer compliance
    \item worker beliefs
    \item reporting and applications
\end{itemize}
\end{column}
\begin{column}{0.32\linewidth}
\textbf{Recognition}
\begin{itemize}
    \item marriage
    \item specialization
    \item family labor supply
\end{itemize}
\end{column}
\begin{column}{0.32\linewidth}
\textbf{Benefits}
\begin{itemize}
    \item partner coverage
    \item spousal insurance
    \item dependent coverage
\end{itemize}
\end{column}
\end{columns}
\vspace{0.2cm}
Legal recognition is an input; actual workplace treatment is an outcome to measure.
\end{frame}

\begin{frame}{Transgender and nonbinary frontier}
\begin{itemize}
    \item Trans and nonbinary labor research faces small cells, privacy, changing categories, disclosure, and legal-documentation issues.
    \item Audit designs observe treatment of a signal at the hiring margin.
    \item Administrative linkage observes earnings and employment for measured subgroups over time.
    \item Survey and list experiments observe stigma, support, and perceived climate.
    \item The frontier contribution is often measurement plus identification, not only a new outcome table.
\end{itemize}
\end{frame}

\begin{frame}{Method map}
\scriptsize
\begin{tabular}{p{0.26\linewidth} p{0.30\linewidth} p{0.31\linewidth}}
\toprule
Method & Best observed margin & Core warning \\
\midrule
Correspondence audit & screening and callbacks & not within-job treatment \\
Large-employer audit & firm concentration of barriers & still one stage of the match \\
Admin identity linkage & earnings and employment dynamics & selected measured group \\
Policy/event study & legal recognition or protection timing & law is not climate \\
Survey/list experiment & stigma, support, hidden attitudes & not direct career outcomes unless linked \\
Benefit records & employer compensation design & need worker take-up and valuation \\
\bottomrule
\end{tabular}
\end{frame}

\begin{frame}{Research opportunities}
\scriptsize
\begin{tabular}{p{0.31\linewidth} p{0.54\linewidth}}
\toprule
Opening & Contribution margin \\
\midrule
Hiring to careers & link screening to entry, wages, promotion, retention \\
Hidden harms & measure climate, concealment, harassment, safety, dignity \\
Worker beliefs & separate employer discrimination from anticipated treatment \\
Firm heterogeneity & locate barriers by employer, occupation, manager, benefit regime \\
Trans and nonbinary measurement & distinguish legal category, identity, perception, disclosure \\
Benefits as compensation & study insurance, leave, family recognition, partner coverage \\
Comparative evidence & test which mechanisms travel across law and enforcement \\
Intersectionality & use powered, privacy-safe heterogeneity designs \\
\bottomrule
\end{tabular}
\end{frame}

\begin{frame}{Week 10 lab}
\begin{itemize}
    \item Reproduce: synthetic Tilcsik-style callback gaps from a randomized identity signal.
    \item Diagnose: compare gay identity-signal and transgender hiring-audit designs; name observed and hidden margins.
    \item Transfer: move from hiring audits to marriage-recognition and employer-benefit designs.
    \item Optional memo: administrative transgender earnings, large-employer audits, or benefit design as compensation.
\end{itemize}
\end{frame}

\begin{frame}{Bridge to final synthesis}
\begin{enumerate}
    \item State the labor object.
    \item Define the identity object.
    \item Name the observed margin.
    \item State the identifying variation.
    \item Separate selection, hiring, within-job treatment, law, and climate.
    \item Say which welfare margins remain hidden.
\end{enumerate}
\vspace{0.2cm}
This is the same final-project discipline as the rest of the course, with harder measurement and richer welfare stakes.
\end{frame}

\end{document}
